\section{Method}
\label{sec:method}

\subsection{Problem setup}
\label{subsec:problem_setup}

Let \(y\in\mathbb{R}^{N}\) denote a target-property window, such as clay volume
or porosity, and let \(u\in\mathbb{R}^{p}\) denote the corresponding dense log
window. Sparse core measurements are modeled as
\begin{equation}
  b = M y + \eta,
  \label{eq:measurement_model}
\end{equation}
where \(M\in\mathbb{R}^{m\times N}\) is a known subsampling matrix, \(m<N\), and
\(\eta\) is measurement noise. The objective is to estimate \(y\) in a target
well using both \(u\) and \(b\).

\begin{assumption}[Windowed generative structure]
\label{ass:generator}
Target-property windows are well approximated by the range of a decoder
\(G:\mathbb{R}^{k}\to\mathbb{R}^{N}\), trained on target windows from training
wells, with \(k<N\).
\end{assumption}

\begin{assumption}[Log-conditioned latent predictability]
\label{ass:latent_prior}
The dense log window \(u\) contains enough information to predict a useful
latent prior \(z_0=h(u)\), where \(h\) is fitted on encoded training windows.
\end{assumption}

\subsection{Generative prior}
\label{subsec:generative_prior}

We train an autoencoder on standardized target windows. Its encoder maps
\(y\) to a latent code \(E(y)\), and its decoder maps a code back to a target
window, \(G(z)\). After training, only \(G\) is used in the inverse problem.
This separates the signal prior from the sparse measurement operator: the
decoder defines the admissible set of plausible property windows, whereas
\(M\) specifies which depths are observed in the target well.

\subsection{Conditional latent prior}
\label{subsec:conditional_prior}

The conditional prior is a ridge regression model
\begin{equation}
  z_0 = h(u),
  \label{eq:latent_prior}
\end{equation}
trained to map log windows to encoded target latents. Ridge regression was
chosen because it is stable in low-data regimes and regularizes correlated log
features \citep{hoerl1970ridge}. The method is not tied to ridge regression:
other predictors can be substituted for \(h\), but the reported paper-facing
variant is CLP-CSGM Ridge.

\subsection{Test-time latent refinement}
\label{subsec:latent_refinement}

At validation and test time, the latent code is refined by
\begin{equation}
  \hat z =
  \arg\min_z
  \underbrace{\|M G(z)-b\|_2^2}_{\text{sparse measurement consistency}}
  +
  \lambda
  \underbrace{\|z-z_0(u)\|_2^2}_{\text{conditional latent prior}},
  \qquad
  \hat y = G(\hat z).
  \label{eq:clp_csgm_objective}
\end{equation}
The regularization parameter \(\lambda\) is selected on validation windows for
each measurement ratio and random seed.

\begin{figure}[t]
  \centering
  \resizebox{\linewidth}{!}{%
  \begin{tikzpicture}[
    node distance=7mm and 8mm,
    >=Latex,
    every node/.style={font=\small},
    block/.style={
      draw,
      rounded corners=2pt,
      align=center,
      minimum width=2.45cm,
      minimum height=0.95cm,
      inner sep=5pt
    },
    io/.style={block, fill=blue!8},
    model/.style={block, fill=green!10},
    latent/.style={block, fill=orange!12},
    solve/.style={block, fill=red!10},
    artifact/.style={block, fill=green!6},
    stage/.style={
      draw,
      rounded corners=4pt,
      dashed,
      inner sep=8pt
    },
    arrow/.style={->, thick},
    guide/.style={->, thick, dashed}
  ]
    % Offline stage
    \node[io] (y) {Training target\\windows \(y\)};
    \node[model, right=of y] (ae) {Train autoencoder\\\(E, G\)};
    \node[latent, right=of ae] (ey) {Encode targets\\\(E(y)\)};
    \node[artifact, right=of ey] (g) {Frozen decoder\\\(G\)};

    \node[io, below=of y] (u) {Training dense\\logs \(u\)};
    \node[model, right=of u] (hfit) {Fit latent prior\\from \((u, E(y))\)};
    \node[artifact, right=of hfit] (h) {Trained prior\\\(h:u\mapsto z_0\)};

    \draw[arrow] (y) -- (ae);
    \draw[arrow] (ae) -- (ey);
    \draw[arrow] (ey) -- (g);
    \draw[arrow] (u) -- (hfit);
    \draw[arrow] (ey) -- (hfit);
    \draw[arrow] (hfit) -- (h);

    \node[stage, fit=(y)(ae)(ey)(g)(u)(hfit)(h),
      label={[font=\bfseries\small]above:Offline training}] (offline) {};

    % Target-well stage
    \node[io, below=20mm of u] (uast) {Target-well dense\\logs \(u_\ast\)};
    \node[latent, right=of uast] (z0) {Latent anchor\\\(z_0=h(u_\ast)\)};
    \node[solve, right=of z0, minimum width=4.4cm] (opt) {Latent refinement\\$\min_z\, \|MG(z)-b_\ast\|_2^2 + \lambda\|z-z_0\|_2^2$};
    \node[model, right=of opt] (yhat) {Reconstruction\\\(\hat y = G(\hat z)\)};
    \node[io, below=of opt] (bast) {Sparse target\\measurements \(b_\ast\)};

    \draw[arrow] (uast) -- (z0);
    \draw[arrow] (z0) -- (opt);
    \draw[arrow] (opt) -- (yhat);
    \draw[arrow] (bast) -- (opt);
    \draw[guide] (h.south) -- node[right] {apply \(h\)} (z0.north);
    \draw[guide] (g.south) -- node[right] {use \(G\)} (opt.north);

    \node[stage, fit=(uast)(z0)(opt)(yhat)(bast),
      label={[font=\bfseries\small]above:Target-well inference}] (online) {};
  \end{tikzpicture}
  }
  \caption{CLP-CSGM workflow. Offline, target windows train an autoencoder and
  dense logs are paired with encoded latents to fit the conditional prior
  \(h\). For a target well, dense logs produce the latent anchor \(z_0\), sparse
  measurements enforce data consistency, and latent optimization is decoded by
  \(G\) to obtain the final estimate.}
  \label{fig:clp_csgm_workflow}
\end{figure}

Figure~\ref{fig:clp_csgm_workflow} summarizes the complete CLP-CSGM workflow and
separates the operations that are learned offline from the operations performed
for a target well. In the offline stage, target windows \(y\) train an
autoencoder, whose decoder \(G\) is retained as the generative range used by the
inverse problem. The same trained encoder maps target windows to latent codes
\(E(y)\); these latents, paired with dense log windows \(u\), define the
supervised training set for the latent prior \(h:u\mapsto z_0\). In the
target-well stage, the dense logs \(u_\ast\) are used only to compute the latent
anchor \(z_0=h(u_\ast)\), while sparse observations \(b_\ast\) enter only
through the data-fidelity term. The final estimate is therefore not a direct
regression from \([u_\ast,b_\ast]\) to \(y_\ast\); it is the decoded output of a
latent code selected to remain close to the log-conditioned prior and consistent
with the measured target entries.


\begin{proposition}[MAP and Tikhonov interpretation]
\label{prop:map}
Assume \(b\mid z\sim\mathcal{N}(M G(z),\sigma_b^2 I)\) and
\(z\mid u\sim\mathcal{N}(z_0(u),\sigma_z^2 I)\). Then the maximum a posteriori
estimate of \(z\) minimizes Eq.~\eqref{eq:clp_csgm_objective} with
\(\lambda=\sigma_b^2/\sigma_z^2\), up to a positive multiplicative constant.
\end{proposition}

\begin{proof}
The negative log posterior is, up to constants independent of \(z\),
\[
  \frac{1}{2\sigma_b^2}\|M G(z)-b\|_2^2
  +
  \frac{1}{2\sigma_z^2}\|z-z_0(u)\|_2^2.
\]
Multiplying by \(2\sigma_b^2\) gives Eq.~\eqref{eq:clp_csgm_objective}.
\end{proof}

\begin{proposition}[Relation to classical CSGM]
\label{prop:csgm_relation}
If \(\lambda=0\), CLP-CSGM reduces to the standard CSGM measurement-consistency
problem \(\min_z \|M G(z)-b\|_2^2\). If \(\lambda\to\infty\), the solution is
forced toward the direct latent prediction \(z_0(u)\), and the method approaches
the decoder-only estimate \(G(z_0(u))\).
\end{proposition}

\begin{proposition}[Existence of a minimizer]
\label{prop:existence}
Assume that the decoder \(G:\mathbb{R}^k \to \mathbb{R}^n\) is continuous and
that \(\lambda>0\). For fixed \(M\), \(b\), and \(z_0\), the objective
\[
  \mathcal{J}(z)
  =
  \|M G(z)-b\|_2^2
  +
  \lambda\|z-z_0\|_2^2
\]
admits at least one global minimizer in \(\mathbb{R}^k\).
\end{proposition}

\begin{proof}
Since \(G\) is continuous and \(M\) is linear, the map
\[
  z \mapsto \|M G(z)-b\|_2^2
\]
is continuous. The quadratic regularization term
\[
  z \mapsto \lambda\|z-z_0\|_2^2
\]
is also continuous and, because \(\lambda>0\), coercive in \(z\). Therefore,
\[
  \mathcal{J}(z)\to+\infty
  \qquad \text{as} \qquad
  \|z\|_2\to\infty .
\]
Hence \(\mathcal{J}\) is continuous and coercive on \(\mathbb{R}^k\), so it
attains its global minimum at some \(\hat z\in\mathbb{R}^k\).
\end{proof}

\begin{remark}[Role of the regularization parameter]
\label{rem:lambda_role}
The parameter \(\lambda\) balances adherence to the log-conditioned latent
anchor against sparse-measurement consistency. Small values of \(\lambda\)
recover behavior closer to measurement-driven generative inversion, while large
values increasingly restrict the solution toward the decoder-only conditional
estimate \(G(h(u))\). Thus, \(\lambda\) does not merely scale a penalty term; it
determines how strongly dense contextual information is allowed to shape the
final reconstruction relative to sparse direct observations. The validation grid
implements this trade-off empirically.
\end{remark}

\begin{remark}[Representation error and decoder mismatch]
\label{rem:representation_error}
The proposed reconstruction is restricted to the image of the decoder \(G\).
Therefore, even when the sparse measurements are noiseless and the optimization
is solved exactly, the reconstruction error cannot in general be smaller than
the mismatch between the true signal and the decoder range. More precisely, if
\(y^\star\) denotes the unknown target signal and
\[
  \mathcal{M}=\operatorname{Im}(G)
\]
denotes the learned manifold, then the method operates under an implicit
approximation error of the form
\[
  \inf_{\tilde y\in\mathcal{M}}\|y^\star-\tilde y\| .
\]
This term is not specific to the conditional latent prior; it is inherited from
the use of a learned generative range itself. Consequently, CLP-CSGM should be
interpreted as combining two sources of inductive bias: a manifold constraint
induced by the decoder and a conditional latent regularization induced by the
logs.
\end{remark}

\begin{remark}[Interpretation of the ablation study]
\label{rem:ablation_interpretation}
The ablation study separates the contributions of three structural components:
the learned manifold \(G\), the conditional latent anchor \(z_0=h(u)\), and the
sparse-measurement consistency term. The prior-only variant isolates the
conditional latent prior without test-time assimilation, whereas the
measurement-only variant isolates the generative inverse problem without dense
contextual anchoring. The complete CLP-CSGM objective combines both mechanisms.
Thus, improvements of the full method over both ablated variants empirically
support the claim that its gains arise from the interaction between contextual
latent regularization and sparse-measurement refinement, rather than from either
component alone.
\end{remark}

\begin{remark}[Theoretical scope]
CSGM theory provides recovery results when the measurement operator satisfies
conditions such as S-REC over the generator range \citep{bora2017compressed}.
Our subsampling operators and learned autoencoder are practical choices for
well-log windows, not a setting in which those guarantees are proven directly.
The theory is used here to motivate the generative-prior formulation; empirical
validation establishes performance for the studied petrophysical tasks.
\end{remark}

\subsection{Limitations}
\label{subsec:method_limitations}

The method inherits representation error from the decoder: if a target profile
is outside the learned range, no latent optimization can recover it exactly.
The optimization problem is nonconvex, so restarts and validation are important.
Finally, the method is computationally heavier than direct regression because it
solves a latent inverse problem for each validation and test window.
